% RESEARCH_DRAFT_SHA256: db2579113fdeb2c516eff8d60029cbdb88863094306a9d220179fb12dd481628
% Layout and source-reference wrappers expanded for the web renderer.
\documentclass[11pt]{article}
\usepackage[margin=25mm]{geometry}
\usepackage{amsthm}
\usepackage{fontspec,unicode-math}
\setmainfont{Latin Modern Roman}
\setmathfont{Latin Modern Math}
\usepackage{amsmath,amssymb,mathtools}
\usepackage{xurl,hyperref}
\hypersetup{colorlinks=true,urlcolor=blue,linkcolor=blue}
\setcounter{secnumdepth}{0}
\setlength{\parindent}{0pt}
\setlength{\parskip}{5pt}
\setlength{\emergencystretch}{3em}
\newcommand{\sref}[2]{\hyperlink{source:#1:#2}{[S#2]}}
\newcommand{\cataloguescope}[1]{\paragraph{Source question.}#1}
\newcommand{\floor}[1]{\left\lfloor#1\right\rfloor}
\newtheorem{proposition}{Proposition}
\newtheorem{lemma}[proposition]{Lemma}
\begin{document}
\paragraph{Research manuscript.} OpenAI \texttt{gpt-6-astra}, reasoning \texttt{ultra}; 11 October 2026.
\paragraph{Preserved source body.} The complete catalogue mathematical text and sources follow with layout and source links expanded, including source qualifications and its historical preparation notes. New research begins at the section headed Research attempt.

\section{Subquadratic Approximation for Longest Common Subsequences}\label{statement}

\subsection{Definitions and mathematical statement}

\paragraph{Shared notation.}Standard number systems are denoted by $\mathbb N,\mathbb Z,\mathbb Q,\mathbb R,\mathbb C$; index ranges, computational models and specialized notation follow the cited problem source and the definitions below.


A common subsequence deletes symbols from each string while preserving order; longest common subsequence maximizes the retained length. An approximation for this maximization problem must preserve the stated fraction of its optimum.
\paragraph{Statement recorded in the supplied source.}
Does there exist, for some $\varepsilon>0$, a $1.1$-approximation algorithm for the longest common subsequence of two length-$n$ strings that runs in time $O(n^{2-\varepsilon})$ for every alphabet size? \href{https://doi.org/10.4171/owr/2018/51}{[S2]}
\paragraph{Scope qualification.}
The meaning of factor $1.1$, alphabet encoding, and deterministic/randomized algorithm model are not specified in the short question.
\paragraph{Source status.}
The supplied catalogue labels this record partially solved. Its source-status metadata was checked against the supplied snapshot on 2026-10-11. This import does not independently establish current open status or mathematical correctness. The catalogue literature assessment is dated 2026-08-21.
\subsection{Short English statement}
Approximate longest common subsequences within the stated factor in truly subquadratic time for every alphabet.
\subsection{Sources}
\begin{itemize}
\item[S1] Ulam AI and the UnsolvedMath Contributors, UnsolvedMath: A Curated Collection of Open Mathematics Problems, catalogue entry. Catalogue curation and metadata: CC BY 4.0. The supplied source question is retained; definitions and the English summary are editorial additions.
\url{https://huggingface.co/datasets/ulamai/UnsolvedMath}.
\item[S2] Complexity Theory (2019), Oberwolfach Reports, DOI 10.4171/owr/2018/51, pp. 3025--3080
\url{https://doi.org/10.4171/owr/2018/51}.
\item[S3] New Frameworks for Hardness of Approximation
\url{https://ieee-focs.org/FOCS-2017-Papers/3464a025.pdf}.
\item[S4] Approximation Algorithms for LCS and LIS
\url{https://epubs.siam.org/doi/10.1137/20M1316500}.
\item[S5] Improved Approximation for LCS over Small Alphabets
\url{https://arxiv.org/abs/2105.03028}.
\textit{Underlying Oberwolfach Reports material retains its source-specific licence, including CC BY-SA 4.0 where applicable; see THIRD\_PARTY\_NOTICES.md.}
\end{itemize}
\paragraph{Preparation and review.}
Original question attributed to the cited source and UnsolvedMath contributors. Codex (GPT-6; exact serving revision not exposed) assisted with definitions, English summaries, source inspection and formatting. No new human mathematical authorship or independent proof verification is asserted. Catalogue identity checked against the complete supplied registry and existing repository statements; known flagged repetitions and obvious identical conjectures excluded. Full mathematical-equivalence/literature review has not been performed. Statement preparation and final review followed contribution prompts C and D; automated checks verify structure and identity, not open status.
\label{end:statement}

\section{Research attempt}
\subsection{Literature, source context and scope}
The full original catalogue body and source qualifications are preserved above. The relevant primary workshop passages were read in the cached publisher report. Historical claims and subsequent attributed results are distinguished from the mathematical work below. A restricted example, standard lemma or known theorem is not counted as new progress on the entire original target. Missing assumptions and repeated or rhetorical extracts are not silently turned into different solved problems.

\subsection{Speeding up by aligned blocks can destroy every match}
Interpret a factor $1.1$ for this maximization task as producing a common subsequence of length at least $\operatorname{LCS}/1.1$. This is stated for the attempted algorithm, not imposed on every source convention. Divide both length-$n$ strings into the same consecutive blocks of size $b$, where $b\mid n$, solve LCS exactly on each pair of equally positioned blocks, and concatenate these subsequences. Standard dynamic programming on a block costs $O(b^2)$, so the total is $O(nb)$ in a unit-cost symbol-comparison model.
\begin{proposition}
Whenever $n$ is even and $b\mid n/2$, this aligned-block algorithm has an instance with true LCS length $n/2$ and output length zero.
\end{proposition}
\begin{proof}
Use $n$ distinct symbols and set
\[
X=(1,2,\ldots,n),\qquad
Y=(n/2+1,\ldots,n,1,\ldots,n/2).
\]
In $X$ every symbol of the first half precedes every symbol of the second half; in $Y$ the order of these two groups is reversed. A common subsequence therefore cannot use a symbol from each group. It can use all $n/2$ symbols from either group, so the optimum is exactly $n/2$.
Because $b$ divides $n/2$, equally positioned blocks of $X$ and $Y$ lie in different symbol groups and have disjoint alphabets. Their LCS lengths are all zero, and so is their concatenation. The concatenation is always a valid common subsequence for general inputs because block orders agree, but validity supplies no approximation guarantee.
\end{proof}
For example $b$ of order $\sqrt n$ gives truly subquadratic running time for the proposed restriction while this family still has zero approximation ratio. Searching more offsets would repair this particular half-rotation, but the proof above does not show that a bounded collection of offsets handles arbitrary alignments.
\subsection{Verification and remaining gaps}
Exact dynamic programming was run on the full strings and aligned blocks for several divisible pairs $(n,b)$, obtaining the predicted $n/2$ versus 0. The alphabet has size $n$, which is permitted by the all-alphabet target, with symbols encoded in $O(\log n)$ bits. The primary small-alphabet approximation paper was read at abstract level and its improvements are not reproduced here. No $1.1$ approximation or general conditional lower bound is established; only a concrete attempted speedup is disproved.

\section{References and source attribution}
Primary workshop report: \url{https://doi.org/10.4171/owr/2018/51}. Relevant passages read: 2018/51 p. 3060, two-string LCS discussion and Open Question 2.

\url{https://arxiv.org/abs/2105.03028}. Read primary Akmal--Vassilevska Williams abstract giving improved truly subquadratic approximations for fixed alphabets; no inference that this settles the1.1 all-alphabet target.

All retrievals were checked on 11 October 2026. Original question attribution remains in the preserved Sources section. External theorems are attributed and not claimed as new discoveries or fully verified merely by citation.


\paragraph{Integration review.} Codex (GPT-6; exact serving revision not exposed) reviewed the written arguments and expanded layout and source-reference wrappers. Repository metadata and the production web renderer were checked using Python 3.13 and JavaScript with MathJax 3.2.2. This model-assisted review is not independent human proof verification; no PDF compilation or formal proof checking is asserted.

\section{Completion Estimate}
\noindent\textbf{COMPLETION: 0\%}
No substantiated new progress toward the entire original target is claimed. The worked results and failed route are delimited above; a known result, illustrative model, or auxiliary obstruction is not counted as original-target completion.

\section{Final status and attribution}
\textbf{Attempt status: unresolved.} This describes this research attempt, not a new ruling on the source catalogue's status or on external literature claims.
The mathematics and this self-audit were generated by OpenAI \texttt{gpt-6-astra}, reasoning \texttt{ultra}, on 11 October 2026. Codex (GPT-6; exact serving revision not exposed) orchestrated the batch, reviewed the written mathematical arguments, and contributed formatting/repository review. This additional model review is not independent human verification. No human mathematical authorship, independent human proof checking, or formal verification is claimed. The historical preparation paragraphs in the reproduced catalogue body concern that body only.

Contribution prompt A was applied to the complete definition and supplied source record, followed by the adversarial checks of prompt B. Local reproduction files are \texttt{batch022/checks.py}, \texttt{computations.json}, and the per-ID audit JSON. They were not public links at preparation time. Numerical checks are finite and do not prove unrestricted results. No LaTeX engine was found locally; compilation was not performed by this research agent.

\end{document}
